\cotefile{77}
\scanpage{50}
\deleted{Donc si $\Delta'$ est la \unsure{sgn-section|2-sect°} de $\check{\mathbb V}(L)$ \unsure{défini|définie} dans le th, on a $x\in\Gamma\Delta$ ssi $2x\in\Gamma\Delta'$, i.e. $\Delta'=2\Delta$} d'où le résultat.

\uline{Corollaire} \unsure{Tout} \deleted{\unsure{schéma}} module quadratique $\Delta$ sur $S$ est canoniquement muni d'une opération de $\mu_2\;(\simeq(\pm1)_S)$, grâce à la formule précédente. L'opération $x\mapsto-x$ \deleted{\unsure{de}} sur $\Delta$ n'est autre que l'involution canonique de $\Delta$ du § 7.

Décrivons \deleted{\unsure{de suite|explicitement}} \inserted{\unsure{enfin,|encore,}} dans le cas général, $B$ comme quotient de
\[
\underline{\mathrm{Aff}}(F)\simeq\underline{\mathrm{Aff}}\,\check{\mathbb V}(L)\simeq\underline{\mathrm{Sym}}(L^\vee)\simeq\coprod_{n\geq0} L^{\otimes(-n)}
\]
par l'idéal engendré par le sous-module inversible engendré par
\[
(\underset{\substack{\uparrow\\\Gamma(L^2\oplus L\oplus\mathcal O_S)}}{\underline{\unsure{c-b+1|c+b+1}}})\otimes L^{-2}\ \subset\ \underline O_S+L^{-1}+L^{\otimes(-2)}\tag*{$(93)$}
\]
% slanted
\margin{\unsure{remonter|renvois|remarque} au § 2}
\deleted{alors l'involution $\sigma$ \unsure{correspon|correspondante}} ce qui donne ici
\[
B\simeq\underline{\mathrm{Sym}}(L^\vee)/(1-4D)L^{\otimes-2},
\]
et notons que l'involution canonique est définie par le passage (\deleted{parmi} les sections de $F\subset\check{\mathbb V}(L)$ satisfaisant $\alpha$), par l'équation
\[
\sigma x=b-x\tag*{$(94)$}
\]
(\unsure{remonte à|variante :|remarque :} § 7) ce qui donne ici $x\mapsto-x$, dans le cas $b=0$, qui est celui du th.
% archive sheet number: 49
\scanpage{51}
9. \uline{Cas $S$ de car} 2 (i.e. $2\,1_S=0$)

Dans ce cas \deleted{$2F\simeq\check{\mathbb V}(L)$}, on a iso can.
\[2F\simeq\check{\mathbb V}(L)\tag*{(95)}\]
ou ce qui revient au $=$, la condition $\mathrm{tr}(e_B)=2$ sur une forme linéaire $\mathrm{tr}$ sur $B,\unsure{e_B}$ muni de $e_B$ \deleted{implique \unsure{$\mathrm{tr}$}} \inserted{signifie}
\[\mathrm{tr}(e_B)=0\tag*{(96)}\]
de sorte que la forme trace $\mathrm{tr}$ (\deleted{\unsure{ou}} ou la section $\underline b$.) est une section de $L$
\[\underline b=\mathrm{tr}\in\Gamma(S,L)\tag*{(97)}\]
canoniquement associée à \deleted{\unsure{tte|la}} structure d'alg. quadratique sur $B$. \deleted{Ceci dit, on a ceci} \inserted{D'ailleurs,} le discriminant est donné par
\[D=\underline b^2\tag*{(98)}\]
car en termes d'un scindage \inserted{$s$} on aura \deleted{\unsure{$b(s)$}}
\[\deleted{\unsure{b_s}}\,b(s)=\underline b\qquad\text{(indépendant de $s$ !)}\tag*{(99)}\]
et
\[D=b^2-4c=b^2=\underline b^2\qquad\text{(car $4=0$)}\]

\uline{NB} Ici l'équation $x^2=D$ s'écrit donc $(x^{\deleted{\unsure{2}}}-\underline b)^2=0$, et donne un sous-schéma de $\check{\mathbb V}(L)$ isom. à $\alpha_2(L)$, et sans grande relation avec $\Delta$\ldots

En concl, \inserted{en plus de $\underline b\in\Gamma(S,L)$,} le torseur $F$ sous $L$, et sa classe
\[f\in H^1(S,L)\tag*{(100)}\]
est un \unsure{nouvel} invariant de la situation. L'involution
% archive sheet number: 50
\scanpage{52}
\deleted{\unsure{de}} \uline{canonique de} $\Delta\inserted{\subset F}$ \uline{n'est autre que la translation par} $\underline b$ (cas particulier de (94) en car. 2). \deleted{Donc, pour $\underline b$ fixé, on peut l'exprimer}

Dire que $\sigma$ dans $B$ est une réflexion (\deleted{pas} seult une pseudoréflexion) signifie ici que $\underline b$ est non nul en chaque fibre\deleted{\unsure{s}} (§ 7), ou encore (\unsure{p.c.q.|p.ex.|par ex.} en vertu de (98)) que $D$ est \deleted{\unsure{inversible|partout}} \inserted{section} inversible, ou enfin que $\Delta$ est étale sur $S$.

Dans ce cas, pour $\underline b$ fixé, les $\Delta$ correspondants (i.e. les $q$ correspondantes \ldots) correspondent biunivoquement aux sections du schéma
\[F/\underline b\]

Or \deleted{si on considère l'hom} on aura une suite exacte de schémas en groupes sur $S$
\diagram{\textup{(101)}\quad$\begin{tikzcd}[ampersand replacement=\&,column sep=small]
0\arrow[r]\&(\mathbb Z/2\mathbb Z)_S\arrow[r,swap,"{\varepsilon\mapsto\varepsilon\underline b}"]\&\check{\mathbb V}(L)\arrow[r,"\beta","{x\mapsto x^2+\underline b x}"']\&\check{\mathbb V}(L^{\otimes2})\arrow[r]\&0
\end{tikzcd}$}
% slanted
\margin{Suite exacte d'Artin-Schreyer (en car. 2)}
de sorte que
\[F/\underline b\simeq F\mathbin{\wedge}_{(\check{\mathbb V}(L),\beta)}\check{\mathbb V}(L^{\otimes2})\overset{\mathrm{déf}}=C\tag*{(102)}\]
s'identifie à un torseur \deleted{\unsure{mult}} sous $\check{\mathbb V}(L^{\otimes2})$, et la donnée de $q$ (ou $\Delta$) compatible avec $\underline b$ s'identifie à la donnée d'une section \inserted{$\sigma$} de ce torseur, dont $F$ est un rev. \uline{étale} quadratique $\Delta$ se récupère alors comme l'image inverse par $\sigma$ de ce rev. étale quadratique.
% archive sheet number: 51
\scanpage{53}
\uline{NB} Pour bien faire, il faudrait vérifier que la structure de torseur sous $L^{\otimes2}$ qu'on vient d'obtenir sur le faisceau des $q$ (ou $\Delta$) compatibles avec $\underline b$, est aussi celle du § 6.

\uline{10} \uline{Cas mixte}. \deleted{A) \unsure{Soit $S$} Cas où}
\diagram{$\begin{tikzcd}[ampersand replacement=\&,row sep=small,column sep=small]
\deleted{S=}\&\deleted{S_2}\arrow[d]\&\deleted{\cup}\&\deleted{S_D}\arrow[d]\\
\&\deleted{\text{pts où }2\text{ inv.}}\&\&\deleted{\text{pts où }D\text{ inv.}}
\end{tikzcd}$}
\deleted{i.e. où $\Delta$ est étale en tous les pts de $S$ de car. 2.}

\deleted{La donnée de $\Delta$ équivaut alors à celle d'un rev. quadratique \deleted{\unsure{étale}} \inserted{\unsure{induit|radical}} de $S_2$\deleted{\unsure{$D$|$\partial$}} et d'un}

Nous nous donnons comme premier invariant dominant la situation le couple $(L,D)$ d'un Module inv. $L$ et de
\[D\in\Gamma(S,L^{\otimes2})\tag*{(103)}\]
\deleted{et l'hypothèse du titre est prise en ce sens :}

Il faut définir un $\Delta$ compatible avec ces données séparément sur $S_D$ et $S_2$, et \deleted{\unsure{donner|donc}} une donnée de recollement sur $S_2\cap S_D=S_{2D}$.

a) Sur $S_2$, on sait (§ 8) que $\Delta$ est déterminé à isom can. près par $(L,D)$, comme le schéma des solutions de l'équation
\[x^2=D\qquad(x\in\Gamma(-,L))\tag*{(104)}\]
% archive sheet number: 52
\scanpage{54}
b) Sur $S_D\deleted{\unsure{=S'}}$, on a un rev. [étale], qui doit prolonger le rev. étale qu'on a déjà sur $S_{2D}\deleted{\unsure{=S'_2}}$

On est donc amené \deleted{(oubliant les pts de $V(D)=S-S_D$)} à traiter ce pb de prolongmt, en se rappellant les données $(L',D')$ sur lesquelles il faut ``\unsure{modeler|recoller}'' le prolongmt, \deleted{\unsure{lorsqu'on}} Un premier invariant important est le torseur $F_{\deleted{\unsure{q}}}$\inserted{\deleted{\unsure{$=F|S'$}}} (noté simplement $F$) sous $\check{\mathbb V}(L)$, trivialisé sur $S'_2$. En fait, on a mieux, savoir une section \inserted{$\underline b$} de $2F$, donc une section $\inserted{b_0=}\frac12\underline b$ de $F|S'_2$ \deleted{\unsure{$\subset V|S$}} dont le produit par $f=2$ se prolonge, et est muni d'un prolongmt. En termes de l'inclusion
\[i:S'_f\longleftarrow S'\qquad(f=2\cdot1_{S'})\]
on peut dire que l'on a
\[\deleted{\unsure{\Gamma(S',}}\ \underline L\subset\tfrac1f\underline L\subset i_*(L|S_f)\]
en supposant $f=2$ section régulière de $\unsure{\mathcal O_S|\mathcal O_{S'}}$, et $F$ est un torseur sous $\underline L$ tel que le torseur correspondant sous $\frac1f\underline L$ soit trivialisé. \inserted{La donnée d'} \unsure{Un|un} tel torseur \deleted{et de} équivaut à la donnée d'une section de $\frac1f\underline L/\underline L\simeq\underline L/f\underline L$, \deleted{\unsure{savoir}} obtenue ainsi : on choisit \inserted{(loc !)} une trivialisation \deleted{locale} $s_0$ de $F$, on écrit $b_0$ comme section de $L|S_f$ de façon précise :
\[b_0=s_{\unsure{*|0}}+\tfrac1f x\qquad(x\in\Gamma(S,L))\]
% archive sheet number: 53
\scanpage{55}
et $x$ est défini modulo $f$ indépendamment de $s$ et on trouve ainsi une section canonique $x_0$ de $L/fL$. En fait, \deleted{$L/fL=$} sur $S'_0=\unsure{\mathrm{Var}(f)|\mathrm{Van}(f)}=V(f\underline O_S)=V(2\underline O_S)$, on trouve ici que $x_0=\underline b_{0}(\Delta|S'_0)\in\Gamma(S'_0,L|S'_0)$ (d'après § 9). Ainsi, le torseur $F$ est déterminé \inserted{(si $2\,1_S$ est $\mathcal O_S$-régulier)} par la seule connaissance de \deleted{\unsure{$B$}} $L$ et de la section
\[b_0\in\Gamma(S_0,L_0)\tag*{(105)}\]

De façon précise, la donnée d'un torseur $F$ sous $\underline L$, et d'une section $\underline b$ de $fF$, équivaut \unsure{exactement|exactemt} ($f$ étant une section $\underline O_S$-régulière de $\underline O_S$) à la donnée d'une section $b_0$ de $L_0$, en lui associant le \inserted{(plutôt)} sous-faisceau Zariski de $\frac1f L_{\unsure{0}}$ formé des sections \inserted{locales} $x$ de $\frac1f L$ telles que
% slanted
\margin{marche pour $L$ faisceau quasi-coh. tel que $f$ soit $L$-réguli\unsure{er|ère}\ldots}
\[fx|S_0=b_0.\]

\deleted{$(F,\underline b)$ étant ainsi déterminés par la connaissance de $(L,b_0)$, \inserted{donc $B$ aussi} il reste à déterminer les $q$.}

Dans le cas actuel, on devra avoir
\[b_0^2=D_0\qquad({:}\,D|S_0)\tag*{(106)}\]
ce qui, si $S_0$ est réduit, détermine $b_0$ de façon unique [et l'existence de $b_0$ est une
% archive sheet number: 54
\scanpage{56}
condition nécessaire évidente pour l'existence d'une rev. quadr. $\Delta$ de $S$ correspondant à la donnée $(L,D)$], et on trouve alors, en \deleted{\unsure{fct|fonctions}} \inserted{termes} \inserted{de $(L,b_0)$, le couple} $(F,\underline b\in\Gamma(S,\underset{\scriptscriptstyle\Vert\,2}{f}F))$, donc $(B,e_B,\mathrm{tr})$.

Les structures \deleted{d'algèbre} cherchées \inserted{$q$} sur $B$ correspondent donc aux sections $\underline c$ d'un \inserted{certain} torseur $Q(\underline b)$ sous $\check{\mathbb V}(L^{\otimes2})$. \deleted{Mais \unsure{on}} La fonction \deleted{$\unsure{\Theta(}$} $\Theta:q\longmapsto\deleted{\unsure{\Theta(q)}}=D(q)=-\mathrm{discr}(q)$ \inserted{de $Q(\underline b)$ dans $\check{\mathbb V}(L^{\otimes2})$} \deleted{correspondant donc aux} \inserted{est alors une} fonction linéaire affine, qui satisfait
\[\deleted{\unsure{D}}\Theta(q+\underset{\substack{\uparrow\\\Gamma(L^{\otimes2})}}{\varpi})=\Theta(q)-4\varpi\tag*{(107)}\]
% slanted
\margin{\uline{\unsure{remonter|renvois|remarques}} \uline{\unsure{au|aux} § 6}}
i.e. la partie homogène $\Theta_0$ de $\Theta$ est définie par
\[\Theta_0(\varpi)=-4\varpi\qquad\text{i.e.}\quad\Theta_0=-4\,\mathrm{id}\tag*{(108)}\]

La section $q$ cherchée, doit donc satisfaire la condition
\[\Theta(q)=D\tag*{(109)}\]
ce qui, en termes d'une trivialisation locale $q_0$ de $Q(\underline b)$, permettant d'écrire $q$ comme
\[q=q_0+\varpi\qquad\varpi\in\Gamma(S,L^{\otimes2})\tag*{(110)}\]
s'écrit comme
\[\Theta(q_0)\unsure{-|+|\dot{-}}4\varpi=D\qquad\text{i.e.}\quad4\varpi=\Theta(q_0)-D)\tag*{(111)}\]
% archive sheet number: 55
\scanpage{57}
Comme on a supposé 2 régulier dans $S$, ceci \deleted{\unsure{impliq|implique}} détermine $\varpi$ de façon unique (d'où \uline{unicité} \deleted{de} $q$). Pour la condition d'existence, posons
\[S_1=\unsure{\mathrm{Var}|\mathrm{Van}|V}(4\underline O_S)\subset S\tag*{(112)}\]
(premier voisinage infinitésimal de $S_0$), et considérons. $\Theta(q_0)-D\,|\,S_1=a_1\deleted{\unsure{(q)}}$, il ne dépend pas du choix de $q_0$, et est de lui \unsure{globalt|globalemt|globale} canonique [défini modulo l'existence -- et choix, le cas échéant -- de $b_0$ satisfaisant (106)]
\[a_1(D)=a_1\in\Gamma(S_1,L^{\otimes2}),\tag*{(113)}\]
et la condition \deleted{est que} d'existence de $q$ est
\[a_1(D)=0\tag*{(114)}\]

On a prouvé :

\uline{Th} Soit $S$ schéma tel que $2\cdot1_S$ soit $S$-régulier, $S_0=V(2\cdot\underline O_S)$, $S_1=V(4\underline O_S)$ son premier voisinage infinitésimal. \deleted{On se donne un \unsure{M.} inv. \unsure{soit,} une section $D$ de $L$} \inserted{Alors la cat. des \deleted{ext.} \inserted{rev.} quadratiques} de $S$ (ou des ext. quadratiques de $\mathcal O_S$) équivaut à celle des triples $(L,D,b_0)$ d'un module inversible $L$, \deleted{et} d'une section $D$ de $L^{\otimes2}$ \deleted{satisfaisant} \inserted{et de $b_0\in\Gamma(S_0,L|S_0)$, satisfaisant} les 2 conditions suivantes

a) \deleted{$\exists\,b_0\in\Gamma(S_0,L|S_0)$ telle que} $b_0^2=D_0\overset{\text{déf}}=D|S_0$

b) La section $a_1(D)\in\Gamma(S_1,L^{\otimes2}|S_1)$ (définie plus haut \deleted{est nulle} en termes de $D_1=D|S_1$) est \uline{nulle}.
% archive sheet number: 56
\scanpage{58}
(NB Lorsque $S_0$ est réduit, $b_0$ est défini de façon unique par $(L,D)$, et a) peut s'écrire comme une condition d'\uline{existence} d'une racine carrée de $D_0$\unsure{..|...})

\uline{Corollaire 1} Supposons que l'on ait $(L,D,b_0)$ satisfaisant a), et soit \deleted{\unsure{$T\to S$}} \inserted{$i:U\subset S$ ouvert rétrocompact} tel que \deleted{\unsure{$\forall$ pt $s\in S$} \unsure{$\exists$ un voisinage $U$ de $s$ et $b_0\in\mathcal O_{S_0,s}$}} \deleted{\unsure{b)}} $\exists\,\Delta_U$ sur $U$ \deleted{satisfaisant} \inserted{correspondant} aux \deleted{conditions} \inserted{données} $L|U,D|U,b_0|U$, [i.e on a $a_1(D_{\unsure{0|1}})|(S_1\cap U)=0$].

Alors, si $U\cap S_0$ est schématiquement dense dans $S_0$ [p.ex. $S_0$ réduit et $U\cap S_0$ dense dans $S_0$] [donc $U\cap S_{\unsure{1|2}}$ schématiquem$^{\mathrm t}$ dense dans $S_1$] \deleted{alors} $\exists\,\Delta$ sur $\unsure{S|B}$ correspondant aux données $(L,D,b_0)$.

\uline{Cor 2} \deleted{Soit donné} Supposons $S_0$ réduit et normal en chaque pt, et soit $(L,D)$. \deleted{tel qu'il existe $U$, ouvert rétrocompact de $S$, avec $U\cap S_0$ dense dans $S_0$, et}

Pour qu'il existe une \deleted{$\Delta$ \unsure{cor}} rév. quadratique de $S$ correspondant à $(L,D)$, il f et s. que la condition suivante \deleted{soit} satisfaite : $\forall s$ pt maximal de $S_0$, existe une solution au Pb en le localisé $\operatorname{Spec}(\mathcal O_{S,s})$ \ldots.
% archive sheet number: 57
\scanpage{59}
\unsure{10|18}. Cas mixte.

\uline{Th}. \deleted{\unsure{Soient}} Soit $S$ un schéma tel que $2\cdot1_S$ soit $\underline O_S$-régulier, \deleted{soit $S_0$} \inserted{\deleted{soit} $S_0=\deleted{\unsure{S}}V(2\underline O_S)$ $S_1=V(4\cdot\underline O_S)$.}

Pour tout rev. quadratique $\Delta$ de $S$, considérons le triple $(L,D,b_0)$ formé du module inversible $L\simeq\underline{\det}(\underline{\mathrm{Aff}}(\Delta/S))^{-1}$, de $\deleted{\unsure{\Theta|\partial}}=D(\Delta/S)\in\Gamma(S,L^{\otimes2})$, et de la section $b_0=b(\Delta_0/S_0)\in\Gamma(S_0,L_0)$, [où $\Delta_0=\Delta_{S_0},L=L_{S_0}$], de sorte que l'on a
\[b_0^2=D_0\tag*{(8.1)}\]

a) Le foncteur (\deleted{de} \inserted{entre} groupoïdes)
\[\Delta\longmapsto(L,D,b_0)\]
est pleinement fidèle.

b) \deleted{Supposons} Pour que \deleted{\unsure{tel}} $(L,D,b_0)$ appartienne à l'image essentielle du foncteur précédent, il \deleted{\unsure{suffit} et il faut} \inserted{faut, et il suffit} si $H^1(S,L)\longrightarrow H^1(S\setminus S_0,L)$ est injectif \inserted{(p.ex. si $S$ est affine\ldots)} \deleted{que $D_1=D|S_1$ soit de la forme $b_1^2$, où $b_1\in\Gamma(S_1,L_1)$ $(L_1=L_{S_1})$, avec $b_1|S_0=b_0$, [\unsure{i.e. $D_1=?$ $b_0$}] ou encore} \inserted{qu'on ait $D=b^2-4c$ avec $b\in\Gamma(S,L)$, $c\in\Gamma(S,L^{\otimes2})$.}

c) En tous les cas, pour que $(L,D,b_0)$ appartienne à l'image essentielle, il f et s que l'on ait
\[D_1=\gamma(b_0),\]
où
\[\gamma:\deleted{\unsure{L|S_0}}\ \mathcal O_{S_0}\longrightarrow\underline O_{S_1}\]
est le morphisme de faisceaux \inserted{multiplicatifs} déduit de l'endom.
\[\underline O_{S_1}\longrightarrow\underline O_{S_1},\quad\beta_1\longmapsto\beta_1^2.\]
[compte tenu que $\beta_1\equiv\beta'_1\ (2)\implies\beta_1^2\equiv{\beta'_1}^2\ (4)$].

\deleted{\uline{Dém} a) \unsure{OPS|DPS} \inserted{\unsure{sur|de}} $S$ affine. On note que \inserted{le torseur} $F$ \inserted{sous $L$}, en tant que faisceau \inserted{le prétors\unsure{eur}} Zariski, \deleted{\illegible{1}} \inserted{$F\subset\frac12 L$} est déterminé par $b_0$, comme le sous-faisceau \inserted{de} $\frac12 \deleted{\unsure{L}}\supset\underline L$ image}
% archive sheet number: 58
\scanpage{60}
\deleted{inverse de la section $b'_0$ de $(\frac12 L)/L$ dont l'image par l'isom can. (induit par $2\,\mathrm{id}_{\frac12 L}$) $(\frac12 L)/L\longrightarrow L/2L$ est $b_0$, \deleted{Alors le torseur} \inserted{\deleted{$\underline L$}} \deleted{$2F$ est le sous faisceau de $\underline L$ engendré par et on a}}

Pour a), \unsure{OPS|DPS} $S$ affine, donc les \inserted{rev. quadratiques} $\Delta$ \deleted{envisagés} de $S$ correspondent à des $\underline L$ torseurs $F$ triviaux, donc $\Delta$ \inserted{isom à un schéma} de la forme
\[\Delta(b,c)\subset\underline L,\quad\Delta(b,c)=\{x\mid x^2-bx+c=0\}\]
où $b\in\Gamma(S,\underline L)$, $c\in\Gamma(S,\underline L^{\otimes2})$. Considérons (\unsure{pour|donc} $\underline L,D,b_0$ fixés) la cat. des $\Delta$ correspondants, \deleted{\unsure{Il}} équivalente à celle des $\Delta(b,c)$, avec $b|S_0=b_0$, $b^2-4c=D$, il faut montrer que cette catégorie, \deleted{\unsure{si}} \inserted{qnd} elle est non vide, est équiv. à la cat. ponctuelle, i.e. \inserted{entre} \deleted{\unsure{deux}} de ces objets il existe un isom. unique. Or les \deleted{\unsure{obj|dits}} isom \inserted{entre $\Delta(b,c)$ et $\Delta(b',c')$} correspondent aux $\tau\in\Gamma(S,\underline L)$ tels que
\[\left\{\begin{aligned} b'&=b+2\tau\\ c'&=c+\tau b+\tau^2\end{aligned}\right.\]

Mais \deleted{$b'\equiv b\ (2)$ \unsure{i.e. que|i.e.}} \inserted{$b'_{S_0}=b_{S_0}\;(=b_0)$} signifie bien que $b'=b+2\tau$ (où $\tau$ est, du reste, déterminé de façon unique, 2 étant $\underline O_S$-régulier). \deleted{Donc, quitte à} \deleted{On a alors}

Considérons $c'_1=c+\tau b+\tau^{2}$, on \deleted{\unsure{a}} donc
\[b^2-4c={b'}^2-4{c'_1}^{\deleted{\unsure{2}}}=D\]
et comme $b'^{2}-4c'=D$, on trouve donc $4(c'_1-c')=0$
\[c'_1-c'=0\qquad\text{i.e.}\quad c'\overset{\unsure{d}}=c',\]
% archive sheet number: 59
\scanpage{61}
puisque 2 \deleted{\unsure{(donc 4) est}} (donc 4) est $\underline O_S$-régulier. Donc $\Delta(b,c)\simeq\Delta(b',c')$, et l'isom est unique, cqfd.

b) \deleted{L'argument précédent montre plus généralement, sans hyp. $S$ affine, que la catégorie des $\Delta(b,c)$ correspondant à $(L,D,b_0)$ donnés, est \inserted{$\simeq$ cat.} ponctuelle. Or si la condition $D_1=$ Considérons} $H^1(S,L)\unsure{-|\to}H^1(S,L_0)$ injectif implique que les $\Delta$ correspondants à $(L,D,b_0)$ correspondent à une enveloppe affine qui est un $\underline L$-torseur trivial, i.e. admet un scindage. Ces $\Delta$ sont donc isom. à des $\Delta(b,c)$, \deleted{Or $\unsure{\Delta(b,}$} $b\in\Gamma(S,L)$, $c\in\Gamma(S,L^{\otimes2})$. \deleted{Donc} Pour que $\Delta(b,c)$ donne $b_0,D$, il f et s. que
\[b_{S_0}=b_0,\qquad b^2-4c=D.\]

Donc $(L,D,c_0)$ est dans l'image essentielle ssi $\forall\,b,c$ satisfaisant les conditions précédentes. Ces conditions impliquent évidemment que $D_1=b_1^2$, où $b_1\in\Gamma(S_1,L_1)$ est tel que $b_1^{\deleted{\unsure{*|0}}}|S_0=b_0$.

Inversement, s'il existe un tel $b_1$, \deleted{\unsure{il|à}} les $b$ qui le relèvent forment un torseur sous $L$, \deleted{provient d'une section de $L$, grâce à la s. e.} dont la restriction à $S-S_0$ est trivial, donc \deleted{cette décl. cohomologique associée à $0\to\unsure{L}\xrightarrow{4}L\to L_1\to0$} trivial par l'hypothèse \inserted{sur $S,S_0$}, donc ce torseur admet \deleted{$H^0(S,L)\to H^0(S_1,L_1)\to H^1(S,L)\xrightarrow{4}H^1(S,L)$} une section $b$. Comme $b^2-D\in4\underline L^{\otimes2}$ et que \deleted{Or $H^1(S,L)_2=H^1(S\setminus S_0,L)$} \deleted{\unsure{(2}}4) est $\mathcal O_S$-régulier donc $L^{\otimes2}$-régulier, on peut écrire $b^2-D=\deleted{\unsure{-}}4c$, i.e. $b^2-4c=D$, cqfd

c) La condition d'existence de $\Delta$, \inserted{correspondant à $(L,D,b_0)$,} à cause de a) doit être locale, \deleted{\unsure{il faut}} comme la condition $D_1=\gamma(b_0)$.
% archive sheet number: 60
\scanpage{62}
D'après b), ces conditions sont équivalentes dans le cas affine. Donc elles sont équivalentes dans le cas général.

\uline{NB} \deleted{Pour pouvoir conclure des \unsure{ctes|ces} conditions qu'il existe $q,\ c$} \inserted{Supposons que $(L,D,b_0)$ appartienne à l'image essentielle} i.e. que \deleted{la condition} l'on ait \deleted{\unsure{$\exists\,D_1$}}
\[D_1=\gamma(b_0)\]
quand peut on conclure qu'il existe $b_1\in\Gamma(S_1,L_1)$ avec $b_1{}_{S_0}=b_0$, $b_1^2=D_1$ ? \deleted{Or les $b_1$ en question} Bien entendu, les $b_1$ en question sont \uline{exactt} les $b_1$ qui relèvent $b_0$ (grâce à $D_1=\gamma(b_0)$), l'obstruction est donc dans $H^1(S_0,L_0)$. On trouve donc

\uline{Corollaire 1} Pour que $(L,D,b_0)$ soit dans l'image essentielle, il \deleted{\unsure{faut et que}} suffit, et il faut lorsque $H^1(S_0,L_0)=0$ \deleted{qu'il} (p.ex. $S_0$ affine, a fortiori si $S$ affine) qu'il existe $b_1\in\Gamma(S_1,L_1)$ tel que
\[b_1^2=D_1,\]

\uline{Corollaire 2} Supposons $S_0$ réduit. Alors le foncteur $\Delta\longmapsto(L,D)$ est pl. fidèle, et son image essentielle est formée des couples $(L,D)$ satisfaisant les \deleted{deux} conditions

\deleted{a) $D_0=D|S_0$ est de la forme $b_0^2$ $(b_0\in\Gamma(S_0,L_0))$ (NB \deleted{\unsure{alors}} $b_0$ est uniquement déterminée \deleted{\unsure{$f$}}, donc son existence est une condition locale sur $S_0$) b) $\gamma(b_0)=D_1$ \deleted{ou encore} \deleted{ou encore} la section \deleted{\unsure{$\gamma(b_0)$}} \deleted{$[\,\unsure{\alpha|a}(b_0,D_1)\in\Gamma(S_0,L_0)$}}
% archive sheet number: 61
\scanpage{63}
\deleted{$\alpha(D_1)\in\Gamma(S_0,L_0)$ définie par la condition $\beta(\alpha(D_1))=\gamma(b_0)-D_1$ (où $\beta:L_0\xrightarrow{\sim}\operatorname{Ker}(L_1\longrightarrow L_0)$ est l'isom. canonique induit par la multiplication par 2 dans $L_1$)}
\margin{\deleted{$D_{\unsure{0|1}}\equiv b_{\unsure{1|0}}^2\bmod 2$}}

que $D_1=D\unsure{\upharpoonright|\vert}S_1$ \inserted{$\in\Gamma(S_1,L_1^{\otimes2})$} est \uline{localmt} le carré d'une section de $L_1$. \inserted{Il suffit, et il faut} si $H^1(S_0,L_0)=0$, \deleted{cette condition équivaut au} que \deleted{\unsure{$D_1$}} il existe $b_1\in\Gamma(S_1,L_1)$ tel que
\[b_1^2=D_1.\]

Enfin, il suffit, et il faut si $H^1(S,L)\unsure{-|\to}H^1(S\setminus S_0,L)$ est injectif, qu'il existe $b\in\Gamma(S,\underline L)$ et $c\in\Gamma(S,\underline L^{\otimes\unsure{2}})$ tels qu'on ait
\[b^2-4c=D.\]

En effet, la première assertion résulte du fait que si $S_0$ est réduit, la section $D_0=D_{S_0}\deleted{\unsure{|S_0}}$ admet au plus une \uline{section}. Le reste suit de façon évidente.

\uline{NB} La condition que $D_1=b_1^2$ localement se décompose en deux

a) $D_0=b_0^2$ (où $b_0\in\Gamma(S_0,L_0)$ est uniquement déterminé \inserted{i.e. \deleted{\unsure{$D$}} la classe $\unsure{\partial|\delta}(b_0)\in H^1(S_0,\mu_2)$ est nulle !}

b) $D_1=\gamma(b_0)$, ce qu'on peut aussi écrire sous la forme
\[\alpha(D_1)=0\]
où $\alpha(D_1)\in\Gamma(S_0,L_0)$ est défini par la condition
\[\beta(\alpha(D_1))=\gamma(b_0)-D_1\]
où $\beta:\deleted{\unsure{\operatorname{Ker} S_0}}L_0\unsure{\}}\xrightarrow{\sim}\operatorname{Ker}(L_1\longrightarrow L_0)$
% archive sheet number: 62
\scanpage{64}
est l'iso. can. induit par la multiplication par 2 dans $L_1$.

\uline{Corollaire 3} Supposons $S_0$ réduit et normal.

Pour que $(L,D)$ provienne d'un rev. quadratique de $S$, il f. et s. qu'il en soit ainsi aux points \unsure{$\xi$|$g$|\deleted{$g$}} maximaux de $S_0$.
% archive sheet number: 63
\scanpage{65}
discriminant des formes quadratiques

(16 p\unsure{.|{}})
% archive sheet number: 64
